\documentclass[a4paper,11pt]{article}

\usepackage[T1]{fontenc}
\usepackage[margin=2.5cm]{geometry}
\usepackage{booktabs}
\usepackage{listings}
\usepackage{xcolor}
\usepackage{hyperref}

% The manual sets many \verb literals in justified prose, and \verb forbids a
% break inside itself, so several paragraphs have no feasible line breaking without
% extra stretch.
\setlength{\emergencystretch}{3em}

\lstset{
  basicstyle=\ttfamily\small,
  backgroundcolor=\color{gray!10},
  frame=single,
  framerule=0pt,
  xleftmargin=6pt,
  xrightmargin=6pt,
  breaklines=true,
}

\title{\textsf{beamerthemeInsiDark} --- Minimal NU dark Beamer theme\\[4pt]
  \large Version 1.0}
\author{Pietro Brach del Prever\\
  \small\href{mailto:brachdelprever.p@northeastern.edu}{brachdelprever.p@northeastern.edu}}
\date{2026/09/21}

\begin{document}
\maketitle

\begin{abstract}
	\textsf{beamerthemeInsiDark} is a minimal dark Beamer theme for
	Northeastern University / INSI presentations.  It features a midnight-navy
	header strip with a textured background image, the ColorBrewer
	\texttt{Paired} palette for content, Fira Sans typography, and convenience
	commands for section separators, statement slides, closing slides, figure
	takeaways, and event metadata.
\end{abstract}

\tableofcontents

\section{Requirements}

\begin{itemize}
	\item \textsf{FiraSans} and \textsf{FiraMono} ---
	      \texttt{tlmgr install fira}
	\item \textsf{tikz}, \textsf{pgfplots} (with the \texttt{colorbrewer}
	      library), \textsf{graphicx} (part of standard TeX Live / MacTeX)
	\item \textsf{currfile}, \textsf{etoolbox}, \textsf{microtype} (part of
	      standard TeX Live / MacTeX)
	\item \textsf{glossaries} and \textsf{algorithm2e}, only when the matching
	      theme option is given (Section~\ref{sec:options})
	\item \textsf{sansmath}, only for the figure companion package
	      (Section~\ref{sec:figure})
	\item \texttt{assets/bg-dark.png} and \texttt{assets/logo-insi.png} ---
	      bundled with this package
\end{itemize}

The theme finds its own \texttt{assets/} directory automatically via
\textsf{currfile}, so no assets need to be copied into the document project.

\section{Installation}

Run the bundled script to symlink the package directory into the personal
TEXMF tree:

\begin{lstlisting}
./install.sh              # ./install.sh --uninstall to remove
\end{lstlisting}

Or by hand (macOS/MacTeX):

\begin{lstlisting}
ln -s /path/to/beamer-template \
      ~/Library/texmf/tex/latex/beamerthemeInsiDark
\end{lstlisting}

After that, \verb|\usetheme{InsiDark}| resolves from any document on the
system.  Verify with:

\begin{lstlisting}
kpsewhich beamerthemeInsiDark.sty
\end{lstlisting}

\section{Loading the theme}

\begin{lstlisting}
\documentclass[aspectratio=169, 10pt]{beamer}
\usetheme{InsiDark}
\end{lstlisting}

\subsection{Theme options}
\label{sec:options}

\begin{lstlisting}
\usetheme[acronyms, algorithms]{InsiDark}
\end{lstlisting}

Neither option is on by default, because each pulls in a heavy package that a
typical talk does not need.

\begin{description}
	\item[\texttt{acronyms}] Loads \textsf{glossaries} configured for acronyms
	      only (\texttt{acronym, nonumberlist, nopostdot, nomain, nogroupskip}),
	      plus \verb|\glsdisablehyper|.  The theme calls no
	      \verb|\makeglossaries|, so nothing is printed as a list, \verb|\gls|
	      expands in place, and the document needs no \textsf{makeindex} pass.
	      Supply the acronym definitions yourself, in the preamble or through
	      \verb|\input{acronyms.tex}|:
	      \begin{lstlisting}
\newacronym{ran}{RAN}{radio access network}
\end{lstlisting}
	\item[\texttt{algorithms}] Loads \textsf{algorithm2e}
	      (\texttt{ruled, vlined, linesnumbered}) and restyles it for slides:
	      \verb|\small| body, an \texttt{InsiDark} SemiBold caption, and
	      \texttt{themeMid} line numbers.  Use \verb|\begin{algorithm}[H]| inside
	      a frame.
\end{description}

\section{Metadata commands}

Standard Beamer metadata (\verb|\title|, \verb|\subtitle|, \verb|\author|,
\verb|\institute|, \verb|\date|) is used as usual.  The theme adds two
preamble commands:

\begin{description}
	\item[\texttt{\textbackslash event\{}\textit{name}\texttt{\}}] Sets the
	      conference or event name shown on the title page below the institute
	      line.
	\item[\texttt{\textbackslash addlogo\{}\textit{graphics}\texttt{\}}] Stacks a
	      partner mark under the INSI monogram, flush right with it, on the title
	      and closing slides.  Call it once per mark, in the order the marks
	      should appear.  \verb|\nulogosep| (default \texttt{0.5em}) is the gap
	      between stacked logos.
	      \begin{lstlisting}
\addlogo{\includegraphics[height=0.07\paperheight]{logo-institute}}
\end{lstlisting}
\end{description}

\section{Slide commands}

\subsection{Title slide}

The title page must be placed inside a group that sets the full-bleed dark
background canvas:

\begin{lstlisting}
{
  \setbeamertemplate{background}{}
  \setbeamertemplate{background canvas}{\InsiDarkBg}
  \setbeamertemplate{footline}{}
  \begin{frame}[plain]
    \titlepage
  \end{frame}
}
\end{lstlisting}

\subsection{\texttt{\textbackslash sectionframe}}

Inserts a dark full-bleed slide that displays the current section name
centered in white.  Call it immediately after \verb|\section{...}|:

\begin{lstlisting}
\section{Results}
\sectionframe
\end{lstlisting}

The size comes from the \texttt{section title} Beamer font, so a deck with long
section names can dial it back without redefining \verb|\sectionframe|.

\subsection{\texttt{\textbackslash statementframe[title]\{line\}}}

A dark gradient slide carrying one line, right-aligned and large: the
rhetorical beat between sections, a step above a section divider.  The frame
title is optional.

\begin{lstlisting}
\statementframe[Where this leaves us]{The RIC decides in milliseconds
  what the network takes minutes to measure.}
\end{lstlisting}

Three pieces are exposed separately, for a custom frame that needs only part of
it:

\begin{description}
	\item[\texttt{\textbackslash statement[width]\{line\}}] The text block alone
	      (default width \verb|0.68\textwidth|), styled by the
	      \texttt{statement} color and font.
	\item[\texttt{\textbackslash InsiDarkStatementBg}] The gradient-plus-wave
	      background; use it as a \texttt{background canvas}.
	\item[\texttt{\textbackslash InsiDarkWaveStrokes}] The wave's stroke list, as
	      \textit{offset}\texttt{/}\textit{opacity} pairs, where the offset is in
	      centimetres downward from the topmost stroke.  \verb|\renewcommand| it
	      to retune the wave.
\end{description}

\subsection{\texttt{\textbackslash closingframe\{msg\}\{contact\}}}

Inserts a dark closing slide with a large message and a contact line:

\begin{lstlisting}
\closingframe{Thank you}{your@email.edu}
\end{lstlisting}

\subsection{\texttt{\textbackslash InsiDarkBg}}

An \verb|\includegraphics| call that scales the full-bleed background image
to the slide dimensions.  Use it as the argument to
\verb|\setbeamertemplate{background canvas}{...}| for any custom dark slide.

\subsection{\texttt{\textbackslash nuLogos}}

Overlays the INSI monogram at the top-right corner of the current slide using
a TikZ \texttt{remember picture, overlay} node, plus any marks added with
\verb|\addlogo|.  Used automatically on the title and closing slides; can be
called on any custom dark frame.

\subsection{Backup slides}

Put backup material after \verb|\appendix|.  Beamer hides those sections from
the table of contents and records the main deck's frame total there, so slide 3
of a 44-slide talk with six backup slides reports \texttt{3 / 44} rather than
\texttt{3 / 50}.  The footline reads that total from Beamer's own
\verb|\insertmainframenumber|, and labels backup frames
\texttt{Backup~\textit{N}} instead of numbering them out of the main total, so
\texttt{Backup 2} can be called out in questions.

\begin{lstlisting}
\closingframe{Thank you}{your@email.edu}

\appendix
\begin{frame}{Experimental Setup}
  Detail held back for questions.
\end{frame}
\end{lstlisting}

The total travels through the \texttt{.nav} file, so it is correct from the
second pass onward.  A deck with no \verb|\appendix| reports
\verb|\inserttotalframenumber|, which is what Beamer defines
\verb|\insertmainframenumber| to be.

\section{Content}

\subsection{Frame titles and subtitles}

The dark strip at the top of a content slide is drawn by the
\texttt{background} template and the frame title sits on it, so the title
carries no fill of its own.  A frame takes a second argument, typeset after the
title in \texttt{Paired-A} and preceded by a guillemet:

\begin{lstlisting}
\begin{frame}{Throughput}{measured on the testbed}
\end{lstlisting}

Display type is not hyphenated: the title page, the frame title, and the
statement slide forbid both automatic hyphenation and breaking at an explicit
hyphen, so lines may end short instead.

\subsection{Table of contents}

Entries are marked with a \texttt{Paired-B} square rather than a number, since
an outline says what the parts are, not which section is the fifth.  Beamer
hardcodes \verb|\vskip1.5em| between entries, a value tuned for four sections,
so an eight-section deck runs off the frame; the theme redirects it to
\verb|\nutocsep| (default \texttt{0.7em}).  The length applies only with
subsections hidden, which is when Beamer uses that skip.

\subsection{Lists}

A slide carries figures and short lines, not paragraphs, so the article-tuned
list padding is too loose.  \verb|\tightlist| goes just after
\verb|\begin{itemize}| or \verb|\begin{enumerate}| and tightens it;
\verb|\nulistsep| (default \texttt{2pt}) is the \verb|\itemsep| it applies.
The gap above and below a list is out of its reach, because \verb|\list| spends
\verb|\topsep| and \verb|\parsep| before the first body token is read.
Bullets are plain black, and the three levels are a bullet, an en dash, and a
smaller bullet.

\begin{lstlisting}
\begin{itemize}\tightlist
  \item First key point.
\end{itemize}
\end{lstlisting}

\subsection{Blocks}

Three block environments, each a rounded title bar over a separately rounded
body bar, with the title centered over the body.  An empty title drops the
title bar instead of leaving a blank coloured strip, which gives a callout that
needs the box but no label.

\begin{lstlisting}
\begin{block}{Definition}      % nuGold title bar
  Formal statement.
\end{block}
\begin{alertblock}{Key Result} % Paired-E title bar
  Main contribution.
\end{alertblock}
\begin{exampleblock}{Example}  % Paired-C title bar
  Concrete illustration.
\end{exampleblock}
\begin{block}{}
  Body-only callout: no title bar.
\end{block}
\end{lstlisting}

Each body is a 25\,\% tint of its title colour.  \verb|\nublockpad| (default
\texttt{0.5em}) scales the inset between a bar's coloured edge and its text;
the body bar takes twice the title bar's share of it.

\subsection{\texttt{narrow}}

\verb|\begin{narrow}[<width>]| (default \verb|0.7\textwidth|) centers a
width-limited \texttt{minipage}.  A \texttt{minipage} sets \verb|\textwidth|,
and Beamer sizes blocks, figures, and tables from \verb|\textwidth|, so
anything inside narrows to match.

\begin{lstlisting}
\begin{narrow}[0.5\textwidth]
  \begin{block}{Narrow}
    Sized from the minipage, not the slide.
  \end{block}
\end{narrow}
\end{lstlisting}

\subsection{Figure captions and takeaways}

Two ranks of text sit under a figure.  The caption says what you are looking
at; it belongs to the float, so it stays a real \verb|\caption|.  The takeaway
says why it matters; it is the argument of the slide, it appears under tables
and lists too, and a float has only one caption slot, so
\verb|\takeaway{...}| is a free-standing box outside the figure.

\begin{lstlisting}
\begin{figure}
  \centering
  \includegraphics[width=0.6\textwidth]{plot}
  \caption{Goodput against offered load.}
\end{figure}
\takeaway{The baseline saturates at 30 Mbps; the proposed scheme does not.}
\end{lstlisting}

A slide has no list of figures and no cross-references, so the theme prints the
caption text alone, with no \texttt{Figure~\textit{N}:} label and no separator.
Both ranks are set in the same size and weight, so the pair reads as one
family; colour is what separates them.  A caption is \texttt{themeMid} grey and
describes the figure, a takeaway is \texttt{Paired-B} and argues from it.

\subsection{Font weights}

The body default is Fira Sans Light, and Fira carries several weights between
it and bold.  NFSS has no series key for the middle ones, so the theme exposes
the two worth reaching for, plus SemiBold:

\begin{center}
	\begin{tabular}{lll}
		\toprule
		Declaration            & Text command              & Weight              \\
		\midrule
		\verb|\regularseries|  & \verb|\textregular{...}|  & Regular, one step up from Light \\
		\verb|\mediumseries|   & \verb|\textmedium{...}|   & Medium, two steps up            \\
		\verb|\sbseries|       & \verb|\textsb{...}|       & SemiBold, one notch under bold   \\
		\bottomrule
	\end{tabular}
\end{center}

\verb|\mdseries| is not among them: in NFSS it means series \texttt{m}, which
the \texttt{light} option of \textsf{FiraSans} maps back to Light.

\subsection{Hyperlinks}

Beamer leaves links uncoloured, and \texttt{structure} is near-black, so a URL
or a citation would be indistinguishable from body text.  The theme colours the
outward-facing ones \texttt{Paired-B} and keeps a URL in the body face with
\verb|\urlstyle{same}|.  Internal navigation links take xcolor's \texttt{.},
the current colour, so \textsf{hyperref} repaints nothing and the surrounding
template decides: the outline keeps its \texttt{section in toc} colour, and the
white section head on the dark section page survives.

\subsection{Bibliography}

The theme styles bibliography entries but does not choose an engine, so the
document is free to use \texttt{thebibliography}, BibTeX, or
\textsf{biblatex}.  Items are numbers rather than Beamer's book icon, entries
are \verb|\footnotesize|, and location and note lines are \texttt{themeMid}
grey.  When \textsf{biblatex} is loaded the theme sets \verb|\bibfont| instead,
because \textsf{biblatex} ignores the Beamer font roles.

\section{Color palette}

Content colors come from the ColorBrewer \texttt{Paired} 12-color set, loaded
via the \textsf{pgfplots} \texttt{colorbrewer} library.  They are available as
\texttt{Paired-A} through \texttt{Paired-L} and can be used anywhere via
\verb|\textcolor{Paired-B}{...}|, \verb|\colorbox{Paired-D}{...}|, etc.
Light/dark pairs: A/B (blue), C/D (green), E/F (red), G/H (orange),
I/J (purple), K/L (yellow / brown).

The theme also defines a handful of reserved neutrals and an accent:

\begin{center}
	\renewcommand{\arraystretch}{1.3}
	\begin{tabular}{lll}
		\toprule
		Name              & Hex               & Role                                        \\
		\midrule
		\texttt{InsiDark} & \texttt{\#01091E} & Title / section / header background         \\
		\texttt{themeMid} & \texttt{\#666666} & Footer / caption text                       \\
		\texttt{themeBg}  & \texttt{\#F4F4F4} & Neutral surface                             \\
		\texttt{nuGold}   & \texttt{\#FDBB45} & Accent: title-slide subtitle, default block \\
		\bottomrule
	\end{tabular}
\end{center}

Block title bars are \texttt{nuGold} (default), \texttt{Paired-E} (alert), and
\texttt{Paired-C} (example).  Every bar carries dark text rather than white,
which is what fixes the colour choice: white on gold is 1.7:1 and dark navy on
gold is 11.4:1, black on \texttt{Paired-E} is 10.1:1 and on \texttt{Paired-C}
15.3:1, where white on the strong \texttt{Paired-D} green is only 3.6:1.  So
alert and example titles sit on the light half of their pair, not the strong
one.  Alerted text (\verb|\alert{...}|) is \texttt{Paired-F}.  Itemize and
enumerate bullets are plain black.

\section{Lengths}

Every length is set with \verb|\setlength| in the document preamble.

\begin{center}
	\renewcommand{\arraystretch}{1.3}
	\begin{tabular}{lll}
		\toprule
		Length                & Default            & Controls                                     \\
		\midrule
		\verb|\nuframeht|     & \texttt{3.15em}    & Height of the dark frame-title strip         \\
		\verb|\nutocsep|      & \texttt{0.7em}     & Gap between \verb|\tableofcontents| entries  \\
		\verb|\nulistsep|     & \texttt{2pt}       & \verb|\itemsep| applied by \verb|\tightlist| \\
		\verb|\nublockpad|    & \texttt{0.5em}     & Inset inside a block bar                     \\
		\verb|\nulogosep|     & \texttt{0.5em}     & Gap between stacked logos                    \\
		\bottomrule
	\end{tabular}
\end{center}

The background clip follows \verb|\nuframeht|: the shift that reveals the
bottom strip of the image is computed at begin-document, once
\verb|\paperheight| is known.

\section{Standalone figures}
\label{sec:figure}

\textsf{beamerthemeInsiDark-figure} is the companion package for the
externalized-TikZ model: one PDF per figure, each compiled on its own, which
the deck then \verb|\includegraphics|.  Such a figure is not a Beamer document,
so it cannot \verb|\usetheme{InsiDark}| and falls back to Computer Modern at
\texttt{OT1}.  Two things then break, and the theme causes both: the face stops
matching the slide, and the \texttt{OT1}/\texttt{T1} mismatch misrenders
\texttt{<}, \texttt{>}, \texttt{|} and accented characters, in the figures
only.  Load the package from the figure's preamble, not the deck's:

\begin{lstlisting}
\usepackage{beamerthemeInsiDark-figure}
\end{lstlisting}

It sets \texttt{T1}, Fira Sans Light, and sans math (\textsf{sansmath}), and
nothing else: a figure's data colours and plot styles belong to the document.
\texttt{install.sh} symlinks \texttt{tex/latex/beamerthemeInsiDark} into \texttt{TEXMFHOME}, so
there is no separate install step; on Overleaf, upload the file into the
project alongside the theme.

\section{Full example}

The file \texttt{main.tex} in this package is a complete working template that
demonstrates every slide type and command described above: title slide,
outline, section separators, a statement slide, frame subtitles, font weights,
two columns, all four block forms, \texttt{narrow}, a \textsf{pgfplots} figure
with caption and takeaway, a table, an \textsf{algorithm2e} listing, a
bibliography, the closing slide, and backup slides after \verb|\appendix|.  It
compiles with:

\begin{lstlisting}
latexmk -pdf main.tex
\end{lstlisting}

\section{License}

This work is released under the \LaTeX\ Project Public License,
version~1.3c or later.
See \url{https://www.latex-project.org/lppl.txt}.

\end{document}
